\id{ҒТАМР 52.47.27}{}

\begin{header}
\swa{М.Д.~Бисенгалиев, Р.Т.~Сулейменова, Мухаммад~Насир, Г.Ш.~Досказиева, Ж.К.~Зайдемова, О.Ш.~Тулегенова, Ж.Б.~Шаяхметова, А.С.~Каримов}{ҚАБАТТЫ ГИДРАВЛИКАЛЫҚ ЖАРУ ТЕХНОЛОГИЯСЫН ҚОЛДАНУ ЖӘНЕ ОНЫ ТҮРЛЕНДІРУ}

\tsp{1}М.Д.~Бисенгалиев\envelope,
\tsp{1}Р.Т.~Сулейменова,
\tsp{2}Мухаммад~Насир,
\tsp{1}Г.Ш.~Досказиева,
\tsp{1}Ж.К.~Зайдемова,
\tsp{1}О.Ш.~Тулегенова,
\tsp{1}Ж.Б.~Шаяхметова,
\tsp{1}А.С.~Каримов
\end{header}

\begin{affil}
\tsp{1}С.Өтебаев атындағы Атырау мұнай және газ университеті, Атырау, Қазахстан,

\tsp{2}Химия ғылымдары мектебі Малайзия зерттеу университеті, Пинанг, Малайзия

\corrauthor{Корреспондент-автор: maks\_bisengali@mail.ru}
\end{affil}

Бұл мақалада мұнай ұңғымаларының өнімділігін арттырудың негізгі
әдістерінің бірі ретінде, әсіресе төмен өткізгіш коллекторларды игеру
кезінде қабатты гидравликалық жару (ҚГЖ) қолданудың теориялық және
практикалық негіздері қарастырылған. Қабаттың сүзу-сыйымдылық
қасиеттерін сенімді анықтау қажеттілігіне басты назар аударылады,
өйткені дәл осы сипаттамалар ҚГЖ жүргізудің сәттілігін және оны
қолданудың экономикалық орындылығын анықтайды. Юра коллекторларының
құрылымы мен қасиеттерін, сондай -- ақ өнімділіктің төмендеуінің негізгі
себептерін, соның ішінде колматация мен қабатпен гидродинамикалық
байланыстың нашарлауын ескеретін «ұңғыма -- кенжар маңы аймағы -- қабат»
жүйесінің сүзу моделі салынды.

Талдаудың негізінде айдау және өндіру ұңғымаларында ҚГЖ жүргізудің
нәтижелері аудандық игеру жүйелерінде өндірудің өсуі мұнайдың тұтқырлығы
мен айдалатын судың арақатынасына байланысты екендігі анықталды. Ең
үлкен әсерге ҚГЖ-ды біріктіретін және жақын маңдағы ұңғымалардың кенжар
маңындағы аймақтарын өңдейтін кешенді әсер ету арқылы қол жеткізілетіні
ерекше атап өтіледі. Бұл тәсіл қабатқа біркелкі әсер етеді, дренаж
аймағын қамтуды арттырады және өндірістің жоғары деңгейін ұзақ уақыт
сақтауға ықпал етеді.

Жасалған тұжырымдар кен орнының геологиялық-техникалық параметрлерін
егжей-тегжейлі талдауға негізделген өндіруді қарқындату жөніндегі
іс-шараларды кешенді жоспарлау қажеттілігін растайды.

{\bfseries Түйін сөздер:} гидравликалық қабаттың жарылуы, ұңғымалардың
өнімділігі, мұнай өндіруді арттыру, өндіруді қарқындату, кенжар маңы
аймағы, сүзу-сыйымдылық қасиеттері, кен орындарын игеруді оңтайландыру.

\begin{header}
ПРИМЕНЕНИЯ ТЕХНОЛОГИЙ ГИДРОРАЗРЫВ ПЛАСТА И ЕГО МОДИФИКАЦИЙ

\tsp{1}М.Д.~Бисенгалиев\envelope,
\tsp{1}Р.Т.~Сулейменова,
\tsp{2}Мухаммад~Насир,
\tsp{1}Г.Ш.~Досказиева,
\tsp{1}Ж.К.~Зайдемова,
\tsp{1}О.Ш.~Тулегенова,
\tsp{1}Ж.Б.~Шаяхметова,
\tsp{1}А.С.~Каримова,
\end{header}

\begin{affil}
\tsp{1}Атырауский университет нефти и газа им. Сафи Утебаева, Атырау, Казахстан,

\tsp{2}Школа химических наук Исследовательского Университета Малайзии, Пинанг, Малайзия,

e-mail: maks\_bisengali@mail.ru
\end{affil}

В статье рассмотрены теоретические и практические основы применения
гидравлического разрыва пласта (ГРП) как одного из ключевых методов
повышения продуктивности нефтяных скважин, в особенности при разработке
низкопроницаемых коллекторов. Основное внимание уделено необходимости
достоверного определения фильтрационно-ёмкостных свойств пласта,
поскольку именно эти характеристики во многом определяют успешность
проведения ГРП и экономическую целесообразность его применения.
Построена фильтрационная модель системы «скважина -- призабойная зона --
пласт», учитывающая строение и свойства юрских коллекторов, а также
основные причины снижения продуктивности, включая кольматацию и
ухудшение гидродинамической связи с пластом.

На основе анализа результаты проведения ГРП в нагнетательных и
добывающих скважинах при площадных системах разработки установлено, что
прирост добычи зависит от соотношения вязкостей нефти и закачиваемой
воды. Особо подчёркивается, что наибольший эффект достигается при
комплексном воздействии, совмещающем ГРП и обработки призабойных зон
близлежащих скважин. Такой подход обеспечивает более равномерное
воздействие на пласт, увеличивает охват зоны дренирования и способствует
длительному поддержанию повышенного уровня добычи.

Сделанные выводы подтверждают необходимость комплексного планирования
мероприятий по интенсификации добычи, основанного на детальном анализе
геолого-технических параметров месторождения.

{\bfseries Ключевые слова:} гидравлический разрыв пласта, продуктивность
скважин, повышение нефтеотдачи, интенсификация добычи, призабойная зона,
фильтрационно-емкостные свойства, оптимизация разработки месторождений.

\begin{header}
APPLICATIONS OF HYDRAULIC FRACTURING TECHNOLOGIES AND THEIR MODIFICATIONS

\tsp{1}M.D.~Bissengaliyev\envelope,
\tsp{1}R.Т.~Suleimenova,
\tsp{2}M.~Nasir,
\tsp{1}G.Sh.~Doskazieva,
\tsp{1}Zh.R.~Zaidemova,
\tsp{1}O.Sh.~Tulegenova,
\tsp{1}Zh.B.~Shayakhmetova,
\tsp{1}A.S.~Karimova
\end{header}

\begin{affil}
\tsp{1}S. Utebaev Atyrau University of Oil and Gas, Atyrau, Kazakhstan,
       
\tsp{2}School of Chemical Sciences, Universiti Sains, Penang, Malaysia,

e-mail: maks\_bisengali@mail.ru
\end{affil}

The article discusses the theoretical and practical foundations of
hydraulic fracturing (HF) as one of the key methods for increasing the
productivity of oil wells, especially when developing low-permeability
reservoirs. The main focus is on the need for reliable determination of
the filtration and storage properties of the formation, since these
characteristics largely determine the success of hydraulic fracturing
and the economic feasibility of its application. A filtration model of
the `well -- bottomhole zone - formation' system has been constructed,
taking into account the structure and properties of Jurassic reservoirs,
as well as the main reasons for productivity decline, including clogging
and deterioration of hydrodynamic communication with the formation.

Based on the analysis of the results of hydraulic fracturing in
injection and production wells in area development systems, it has been
established that the increase in production depends on the ratio of the
viscosities of oil and injected water. It is particularly emphasized
that the greatest effect is achieved through a comprehensive approach
combining hydraulic fracturing and treatment of the bottomhole zones of
nearby wells. This approach ensures a more uniform impact on the
formation, increases the coverage of the drainage zone and contributes
to the long-term maintenance of increased production levels.

The conclusions confirm the need for comprehensive planning of measures
to intensify production based on a detailed analysis of the geological
and technical parameters of the field.

{\bfseries Keywords}: hydraulic fracturing, well productivity, oil recovery
enhancement, production intensifica\-tion, bottomhole zone, filtration and
storage properties, field development optimisation.

\begin{multicols}{2}
{\bfseries Кіріспе.} Қабатты гидравликалық жару мұнай мен газ конденсатын
өндірудің ағымдағы деңгейін, айдау ұңғымаларының қабылдау қабілетін
арттыру әдістерінің бірі болып табылады, бұл төмен өткізгіштігі бар
тығыз жыныстарда (құмтастар, әктастар, доломиттер және т.б.) терең
жатқан горизонттарды сынауға мүмкіндік береді. Оның мәні сұйықтықты
жоғары қысыммен айдау кезінде (таудан жоғары) табиғи заттардың ашылуы
және өнімді қабатта жасанды жарықтар пайда болады. Жарықтар артық
қысымды алып тастағаннан кейін олардың жабылуына жол бермейтін бекітетін
агенттерді енгізу арқылы бекітіледі.

Заманауи технология терең енетін қабатты гидравликалық жару (ТҚГЖ)
қабатқа әсер ету аймағын едәуір ұлғайтуды, соның ішінде іргелес
ұңғымалар арасындағы жарықтар жүйесін қосуды қарастырады.

Соңғы жылдары жыртылу нұсқалары кең таралуда: қышқыл, көбік, сілтілер
мен мұнай негізіндегі ерітінділерді қолдану.

Қабатты гидравликалық жару (ҚГЖ) төмен өткізгіш, нашар ағызылатын
коллекторларды ашатын ұңғымалардың өнімділігін арттырудың ең тиімді
әдістерінің бірі болып табылады.

Қабатты гидравликалық жару -сұйықтықты жоғары қысыммен айдау арқылы
өнімді қабатта жарықтар жасау, содан кейін оларды пропантпен бекіту. Бұл
төмен өткізгіш коллекторларды, сондай-ақ үлкен радиусты қабаттың
ластанған кенжар аймақтарын дамытудың тиімді әдістерінің бірі.

Кеуекті кеңістіктегі су мен мұнайдың белгілі бір арақатынасында
өнімділігі төмен коллекторларда әлсіз қозғалмалы жүйе пайда болады, бұл
қабаттардың мұнай өндірісін арттырудың әдеттегі әдістерінің төмен
тиімділігін тудырады. Мұндай аймақтарды жою үшін қысымның айтарлықтай
өзгеруін жасау керек. Бұл мәселені шешудің ең жақсы тәсілдерінің
бірі-ҚГЖ.

{\bfseries Материалдар мен әдістер.} Мұнай кен орындарында ҚГЖ-ды кеңінен
қолдану тәжірибесі көрсеткендей, бұл әдіс кен орнын немесе аумақты игеру
тәсілі ретінде жобалау кезінде бүкіл қабат жүйесін, ұңғымалардың
геометриялық орналасуын және олардың өндіру және айдау ұңғымаларын
өңдеудің әртүрлі комбинацияларындағы өзара әсерін ескере отырып, ең
жоғары тиімділікті қамтамасыз ете алады.

Қабатты гидравликалық жару процесі жасанды жасау және ұңғымаға айдалатын
сұйықтықтың жоғары қысымының әсерінен кенжар аймағындағы жыныстардағы
жарықтарды кеңейту болып табылады. Қысымның жоғарылауымен қабаттың
жыныстарында жаңа жарықтар пайда болады немесе бар жарықтар ашылып,
кеңейеді. Бұл жарықтар жүйесі ұңғыманы кенжардан алынған қабаттың өнімді
бөліктерімен байланыстырады. Қысымды төмендеткеннен кейін жарықтардың
жабылуын болдырмау үшін оларға ұңғымаға құйылған сұйықтыққа қосылатын
ірі дәнді құм енгізіледі. Жарықтар радиусы бірнеше ондаған метрге жетуі
мүмкін.

Кәсіпшілік тәжірибе көрсеткендей, гидравликалық жарудан кейінгі
ұңғымалардың шығыны кейде бірнеше ондаған есе артады. Бұл пайда болған
жарықтар бұрыннан бар жарықтарға қосылатындығын және ұңғымаға ағын бұрын
оқшауланған жоғары өнімді аймақтардан келетіндігін көрсетеді.

Қабатқа сүзілетін сұйықтықпен қабаттың жарылуы кезінде жарықтардың пайда
болу механизмі келесідей: ұңғымада сорғы агрегаттарымен пайда болатын
қысыммен жару сұйықтық, ең алдымен, өткізгіштігі жоғары аймақтарға
сүзіледі. Сонымен қатар, тігінен қабаттар арасында қысым айырмашылығы
пайда болады, өйткені қабаттар туралы өткізгіштігі төмен немесе іс
жүзінде өткізбейтін қысымнан көп қысым бар. Нәтижесінде кейбір күштер
өткізгіш қабаттың төбесі мен табанына әсер ете бастайды, үстіңгі
жыныстар деформацияға ұшырайды және қабаттар туралы шекараларда көлденең
жарықтар пайда болады.

Сүзілмеген сұйықтықтың жарылуы кезінде қабаттың жарылу механизмі қалың
қабырғалы тамырлардың жарылу механизміне ұқсас болады. Бұл жағдайда
пайда болған жарықтар, әдетте, тік немесе көлбеу бағытқа ие. Сүзгіден
өткен сұйықтықтың жарылуы кезінде, әдетте, сүзгіден өтпейтін
сұйықтықтардың жарылуына қарағанда, жарылу қысымы айтарлықтай аз болады,
өйткені соңғы жағдайда жыныстардың жарылу механизмі қалың қабырғалы
ыдыстың жарылу механизміне ұқсас. Жыныспен жанасудың үлкен аймағына
байланысты қабатқа енген сүзгіден өткен сұйықтық оған сүзгіден өтпейтін
сұйықтықпен қабаттың ыдырауы үшін қажет болғаннан әлдеқайда аз қысыммен
жарылу үшін жеткілікті үлкен күш береді. Бұрын қабаттардың жарылу қысымы
тау қысымынан және тау жыныстарының жарылуға уақытша төзімділігінен асып
кетуі керек деп есептелген.

Жарықтардың пайда болу механизміне сүйене отырып, жарылу қысымы көптеген
факторларға байланысты болуы керек: тау қысымының таралу мөлшері мен
сипаты, тау жыныстарының беріктігі мен серпімді қасиеттері,
коллекторлардың өткізгіштігі және жарылған сұйықтықтың физикалық
қасиеттері, кен орнының геологиялық құрылымы, гидравликалық жару
технологиясы және т. б.

Гидравликалық жару қолдану технологиясы, ең алдымен, жарықтардың
геометриясын болжауға және оның параметрлерін оңтайландыруға мүмкіндік
беретін тау жыныстарындағы жарықтардың пайда болуы мен таралу механизмін
білуге негізделген. Жарықшақ түзілу процесін математикалық модельдеу
серпімділік теориясының, мұнай-газ қабаттарының физикасының, сүзудің,
термодинамиканың негізгі заңдарына негізделген. Екі өлшемді жарықшақтың
таралуының жалпы мойындалған алғашқы теориялық моделі жұмыстарда
ұсынылған {[}1-3{]}.

ҚГЖ жүргізудің табыстылығының маңызды факторы жарықтар бойымен
тасымалданатын жарылған сұйықтықтың, пропант немесе құм (құмтас)
түріндегі қатты композиттердің (материалдардың) сапасы болып табылады.
Жару сұйықтығының негізгі мақсаты-жарықшақты ашу үшін қажетті энергияны
ұңғыманың бетінен кенжарға беру және пропантты (құмды) бүкіл жарықшақ
бойымен тасымалдау. ҚГЖ сапасына әсер ететін "жару сұйықтығы және
тасымалдаушы сұйықтық" жүйесінің негізгі қасиеттері:

- "таза" және құрамында пропант (құм) бар сұйықтықтың реологиялық
қасиеттері;

- гидравликалық жару кезінде және пропантты (құмды) жарықшақ бойымен
тасымалдау кезінде оның қабатқа ағып кетуін анықтайтын сұйықтықтың
инфильтрациялық қасиеттері;

- сұйықтықтың пропанттың (құмның) ілулі күйдегі жарықшақтың ұштарына
өтуін қамтамасыз ету қабілеті;

- пропант (құм) бөлшектерінің бір-біріне жабысып, кеуекті кеңістік пен
жарықшақты бітеу мүмкіндігі;

- пропант қаптамасының және қоршаған қабаттың ең аз ластануын
(колматациясын) қамтамасыз ету үшін жарылған сұйықтықты оңай және жылдам
шығару мүмкіндігі;

- үзіліс сұйықтығының технологиямен, мүмкін қоспалармен және қабат
сұйықтықтарымен қамтамасыз етілген әртүрлі қоспалармен үйлесімділігі;

- пропанттың (құмның) физикалық қасиеттері.

Гидравликалық жарылыстың технологиялық сұйықтықтары үлкен ашу және қатты
материалмен (пропантпен) тиімді толтыру арқылы жоғары өткізгіштік
жарықшағын жасау үшін жеткілікті динамикалық тұтқырлыққа ие болуы тиіс;
сұйықтықтың ең аз шығынымен қажетті мөлшердегі жарықшақтарды алу үшін
төмен сүзгіш ағып кетулерге ие болуы; жыныстармен және қабат
сұйықтықтарымен үйлесімділікке ие болуы; жару сұйықтығымен жанасатын
қабат аймағының өткізгіштігінің ең аз төмендеуін қамтамасыз етуі;
құбырлардағы үйкеліс қысымының төмен жоғалуын қамтамасыз ету; өңделетін
қабат үшін жеткілікті жылу тұрақтылығына ие болу; жеткілікті жоғары
ығысу тұрақтылығына ие болу; өңдеуден кейін гидравликалық жару қабатынан
және жарықшақтан оңай шығу; коммерциялық жағдайда дайындау және сақтау
кезінде технологиялық болу; коррозиялық белсенділігі төмен болу;
экологиялық таза және қолдануға қауіпсіз болу; салыстырмалы түрде төмен
құны бар {[}4{]}.

Ашық күйдегі жарықтарды бекіту үшін қолданылатын заманауи материалдар -
пропанттар - келесідей жіктеледі: кварц құмдары, орташа және жоғары
беріктігі бар синтетикалық пропанттар. Жарықтың өткізгіштігіне әсер
ететін пропанттардың физикалық сипаттамаларына беріктік, түйіршіктердің
мөлшері және гранулометриялық құрамы, сапасы (қоспалардың болуы,
қышқылдарда ерігіштігі), түйіршіктердің пішіні (сфералық және
дөңгелектік) және тығыздық сияқты параметрлер жатады {[}5{]}.

Жарықтарды бекіту үшін негізгі және ең көп қолданылатын материал-құм.
Оның тығыздығы шамамен 2,65 г / см\tsp{3}. Құмдар әдетте
қысу кернеуі 40 МПа-дан аспайтын гидравликалық жару кезінде қолданылады.
Тығыздығы 2,7-3,3 г/см\tsp{3} керамикалық пропанттар 69 МПа
дейін қысу кернеуінде қолданылады. Агломерацияланған боксит және
цирконий оксиді сияқты орташа төзімді пропанттар қысу кернеуі 100 МПа-ға
дейін қолданылады, бұл материалдардың тығыздығы 3,2-3,8
г/см\tsp{3} құрайды.

Пропанттар құмнан 2-2,5 есе үлкен жүктемелерге төтеп бере алады, бірақ
олардың тығыздығы құмның тығыздығынан 40-50\% асады, бұл мұндай
бөлшектерді жарықтарға тасымалдауды қиындатады.

Беріктік-бұл қабаттың тереңдігінде жарықшақтың ұзақ өткізгіштігін
қамтамасыз ету үшін белгілі бір қабат жағдайлары үшін пропанттарды
таңдаудағы негізгі өлшем. Терең ұңғымаларда минималды кернеу көлденең
болады. Тереңдікпен минималды көлденең кернеу шамамен 19 МПа/км-ге
артады. Сондықтан әртүрлі тереңдіктер үшін пропанттардың келесі түрлері
қолданылады: кварц құмдары-2500 м дейін; орташа беріктігі бар
пропанттар-3600 м дейін; жоғары беріктігі бар пропанттар-3500 м-ден
жоғары.

Пропант (құм) бөлшектеріне беріктіктен басқа, фракциялық құрамы бойынша
гидравликалық жарылыстар кезінде ерекше талаптар қойылады. Құм
түйіршіктерінің орташа мөлшері әдетте 0,5-0,88 мм болуы керек {[}3{]}.
Бөлшектердің мөлшері ұлғайған сайын жарықтардың гидроөткізгіштігі
артады, ал азайған сайын тасымалдаушы сұйықтықтың тасымалдау қабілеті
артады деп саналады. Қатты фазалық бөлшектердің концентрациясы маңызды
мәнге ие, оның максималды мәні бөлшектердің коалесценциясының
(бір-біріне жабысуының) басталу шамасынан аспауы керек, нәтижесінде
тасымалдаушы сұйықтықтың бөлшектердің жеткізілу ауқымы төмендейді, ал
түсетін бөлшектер тесіктер мен жарықтарды бітеп тастауы мүмкін {[}4{]}.

Көбінесе түйіршіктердің мөлшері 0,85-0,425 мм болатын пропанттар
қолданылады; сирек 1,7-0,85 мм; 1,18-0,85; мм 0,425-0,212 мм. Пропант
дәндерінің дұрыс мөлшерін таңдау факторлардың тұтас кешенімен
анықталады. Түйіршіктер неғұрлым үлкен болса, пропанттың қаптамасы
соғұрлым көп өткізгіштікке ие болады. Алайда, үлкен фракцияның пропантын
қолдану, жоғарыда айтылғандай, оны жарықшақ бойымен тасымалдау кезінде
қосымша қиындықтармен бірге келеді. Пропанттың беріктігі түйіршіктердің
мөлшерінің ұлғаюымен төмендейді, сонымен қатар әлсіз цементтелген
коллекторларда пропантты ұсақ фракцияны қолданған жөн, өйткені ұсақ
бөлшектерді қабаттан шығару арқылы ірі түйіршікті пропанттың қаптамасы
біртіндеп бітеліп, оның өткізгіштігі төмендейді.

Пропант түйіршіктерінің дөңгелектігі мен сфералығына оның қаптамасының
жарықшақтағы тығыздығы, оның сүзуге төзімділігі, сондай-ақ тау қысымының
әсерінен түйіршіктердің бұзылу дәрежесі байланысты. Пропанттың тығыздығы
пропанттың жарықшақ бойымен тасымалдануы мен орналасуын анықтайды.
Тығыздығы жоғары пропанттарды жарықшақ бойымен тасымалдау кезінде
тасымалдаушы сұйықтықта суспензияда ұстау қиын. Жарықшақты жоғары
тығыздықтағы пропантпен толтыруға екі жолмен қол жеткізуге болады:
пропантты жарықшақтың ұзындығы бойынша минималды тұндырумен
тасымалдайтын тұтқырлығы жоғары сұйықтықтарды қолдану немесе оларды
айдау қарқыны жоғарылаған кезде тұтқырлығы төмен сұйықтықтарды қолдану.
Соңғы жылдары шетелдік фирмалар тығыздықтың төмендеуімен сипатталатын
жеңілдетілген пропанттарды шығара бастады.

Орташа және жоғары өткізгіш қабаттардың гидравликалық жарылуы қазіргі
уақытта ұңғымаларды ынталандырудың ең қарқынды дамып келе жатқан
әдістерінің бірі болып табылады. Жоғары өткізгіш қабаттарда, негізгі
фактор жарықшақтың ұзындығы болып табылатын төмен өткізгіштерден
айырмашылығы, ұңғыманың өнімділігінің артуы жарықшақтың еніне
байланысты. Қысқа жарықтар жасау үшін жарықшақтың соңында пропантты
тұндыру технологиясы қолданылады (TSO "tip-screen-out") {[}5{]}, ол
пропантты, ең алдымен, өңдеу кезінде жұмыс сұйықтығындағы
концентрациясын біртіндеп арттыру арқылы жарықшақтың соңына қарай
итеруден тұрады. Пропанттың жарықшақтың соңында тұндырылуы оның
ұзындығының өсуіне жол бермейді. Тасымалдаушы материалды одан әрі айдау
жарықшақтың енінің ұлғаюына әкеледі, ол 2,5 см - ге дейін жетеді, ал
қалыпты ҚГЖ кезінде жарықшақтың ені 2-3 мм құрайды.

Осылайша, бөлшектердің қасиеттері-беріктігі, тығыздығы, мөлшері,
концентрациясы және т.б., сондай-ақ қатты фазаның қабат сұйықтығымен
өзара әрекеттесу сипаты кеуек кеңістігі мен жарықтардың бітелу дәрежесін
анықтайды. Бұл жағдайда әр операцияны дайындауға үнемі назар аударылады.
Мұндай дайындықтың маңызды элементі-бастапқы ақпаратты жинау және
талдау. ҚГЖ дайындау үшін қажетті деректерді үш топқа бөлуге болады
{[}6{]}:

- қабаттың геологиялық-физикалық жағдайлары (өткізгіштігі, кеуектілігі,
қанықтылығы, қабат қысымы, газ-мұнай және су-мұнай байланыстарының
жағдайы, тау жыныстарының петрографиясы);

- жарықтың геометриясы мен бағдарының сипаттамасы (минималды көлденең
кернеу, Юнг модулі, Пуассон коэффициенті, тау жыныстарының сығылуы және
т. б.);

- жару сұйықтығының, тасымалдаушы сұйықтықтың және пропанттың
қасиеттері.

Ақпараттың негізгі көздері геологиялық, геофизикалық және петрофизикалық
зерттеулердің, ядроның зертханалық талдауының, сондай - ақ микро және
мини -- гидравликалық жарылыстар жүргізуден тұратын кәсіптік
эксперименттің деректері болып табылады.

ҚГЖ жүргізудің сәттілігі операцияларды жүргізу үшін объектіні дұрыс
таңдауға, осы жағдайлар үшін оңтайлы үзіліс технологиясын қолдануға,
өңдеуге арналған ұңғымаларды сауатты таңдауға байланысты.

Әрбір нақты жағдайда ҚГЖ жүргізу туралы шешім қабылдау
тау-кен-геологиялық жағдайларды ескере отырып жүзеге асырылады. Әдетте,
әлеуетті объектінің геологиялық-физикалық қасиеттерін талдау кезінде
келесі ерекшеліктер ескеріледі {[}7-11{]}:

1) бұрын құрғатылмаған қабаттар туралы аймақтарды игеруге қосу арқылы
гидравликалық жарудың жоғары тиімділігін қамтамасыз ететін созылу
бойынша қабаттың гетерогенділігі және қалыңдығы бойынша бөлшектену;

2) мұнай тұтқырлығы 5 мПа·с дейін болған кезде әдетте 0,03
мкм\tsp{2}-ден аспауы керек қабаттың өткізгіштігі (жоғары
өткізгіштік қабаттарында жергілікті ҚГЖ тиімді, бұл негізінен кенжар
аймағын өңдеу құралы ретінде айтарлықтай әсер етеді);

3) өнімді қабатты газ немесе сумен қаныққан коллекторлардан бөлетін
литологиялық экрандардың қалыңдығы мен төзімділігі, ол кемінде 4,5-6 м
болуы керек;

4) қабаттың пайда болу тереңдігі, әдетте, 3500 м-ден аспауы керек және
ҚГЖ технологиясына қойылатын талаптарды анықтайды, атап айтқанда,
гидравликалық жарудан кейін ұңғымалардың шығымын едәуір және ұзақ
уақытқа арттыру үшін жеткілікті, сондықтан ҚГЖ жүргізу шығындарының
өтелуін қамтамасыз етеді;

5) алынатын қорлардың өндірісі, әдетте, 30\% - дан аспауы керек.

Егер жарықтар қабатының сүзу сипаттамалары арасындағы тепе-теңдікті
қамтамасыз ететін жарықтар параметрлерін өңдеу және оңтайландыру үшін
ұңғымаларды таңдау объектінің геологиялық-физикалық қасиеттерін,
жарықтар бағытын, су басу жүйесін және ұңғымаларды орналастыруды
анықтайтын қабаттағы кернеулердің таралуын ескере отырып жүзеге
асырылса, ең жоғары тиімділікке қол жеткізуге болады. ҚГЖ жүргізудің
әсері жекелеген ұңғымалардағы жұмыста бірдей көрінбейді, сондықтан
гидравликалық жару салдарынан әр ұңғыманың шығымының өсуін ғана емес,
сонымен қатар ұңғымалардың өзара орналасуының, қабаттың
гетерогенділігінің нақты таралуының, объектінің энергетикалық
мүмкіндіктерінің әсерін де қарастыру қажет.

Юра қабаттарының коллекторларын дамытудың қазіргі проблемаларының
бірі-төменгі қабаттардың төмен сүзу-сыйымдылық қасиеттерімен және төмен
бастапқы мұнаймен қанықтырумен қатар ұңғымалардың өнімділігіне үлкен
әсер ететін ұңғыма аймағындағы физика-химиялық және гидродинамикалық
процестерді басқару. Өкінішке орай, қазіргі уақытта Юра
қабаттары-коллекторлар бойынша зерттеу және кәсіптік ақпараттың
жеткілікті көлемі жоқ, оның негізінде ұңғыма аймағындағы физика-химиялық
құбылыстар мен сүзу процестерін сандық бағалауға, сондай-ақ коллектордың
өткізгіштігінің төмендеуінің көптеген факторларын ескере отырып,
ұңғымалардың өнімділігін қалпына келтірудің және арттырудың оңтайлы
әдістерін жобалауға болады {[}7-10{]}.

Кен орнының юра қабаттары ұңғымаларының әлеуетін бағалау және олардың
өнімділігін арттыру мәселелерін шешу үшін зерттеу кешенін жүргізу арқылы
осы мәселелерді талдау және шешу тұжырымдамасы, сондай - ақ Юра
коллекторларының сүзу -- сыйымдылық сипаттамаларын және ұңғымалардың
өнімділігін төмендетудің негізгі факторларын ескеретін "ұңғыма-кенжар
маңы аймағы-қабат" жүйесінің сүзу моделін құру ұсынылады. Мұндай
модельді құру ұңғымалар мен қабаттарға гидродинамикалық,
кәсіптік-геофизикалық зерттеулер жүргізуге, сондай-ақ қабаттарға жақын
жағдайларда сүзу процестерінің табиғи өзектерінде физикалық модельдеуге
байланысты бірқатар мәселелерді шешуді талап етеді. Бұл, ең алдымен,
ұңғымалардың маңындағы перфорация аралығындағы сұйықтықтарды сүзу
геометриясын зерттеуге (ұңғымаға жақын және шалғай аймақтар), қабаттың
зақымдану аймағының болуы мен тереңдігін анықтауға, өндіруші
ұңғымалардағы қабаттар туралы жұмыс істемейтін жерлерді анықтауға
қатысты. Екіншіден, коллекторлық қабатқа, демек, бұрғылау, технологиялық
өңдеу кезінде әртүрлі факторлардың ұңғымаларының өнімділігіне әсерін
бағалау үшін кешенді зерттеулер жүргізу қажет. Үшіншіден, осы
міндеттерді шешу бойынша зерттеулер жүргізу жекелеген мамандардың
эпизодтық жұмыстарымен емес, бірыңғай, мақсатты, кешенді бағдарлама
шеңберінде орындалуы тиіс. Төмен өткізгіш коллекторлардағы ағынды
қарқындату мәселелерінің әртүрлілігі мен күрделілігі жан-жақты талдауды
және оларды мамандар арасында кеңінен талқылауды қажет етеді.

Ең аз зерттелген мәселе-перфорация аралығындағы сұйықтықтардың сүзілу
сипатын және ұңғымалардың айналасын бағалау. Қазіргі уақытта ұңғымаларды
механикаландырылған пайдалану жағдайында кен орындарын игерудің ағымдағы
сатысында қабаттың және кенжар маңы аймағының сүзу қасиеттері туралы
деректер, сондай-ақ коллектордың зақымдану аймағының болуы мен тереңдігі
және қабаттың қалыңдығы бойынша аралықтар туралы мәліметтер іс жүзінде
жоқ. Сонымен бірге ұңғымаларды гидродинамикалық және
кәсіптік-геофизикалық зерттеу әдістері, сондай-ақ кенжар маңындағы
аймақтың зақымдануының болуы мен тереңдігін бағалауға және қабаттың
қалыңдығы бойынша жұмыс істемейтін аралықтарды анықтауға мүмкіндік
беретін математикалық модельдер бар және тиімді қолданылады. Бұл ақпарат
ұңғымалардың өнімділігін қалпына келтіру және арттыру технологияларын
жоспарлау кезінде шешуші болып табылады.

Кенжар маңы аймағының өткізгіштігін қалпына келтіру және арттыру үшін
кез келген тиімді қарқынды құрамды қолдану, егер қабаттың қалыңдығы,
сондай-ақ зақымдану аймағының болуы мен тереңдігі бойынша сүзу
гетерогенділігі туралы деректер болмаса, көп жағдайда оң нәтижелерге
әкелмейді.

Осы әдістерді, сондай-ақ қазіргі заманғы стационарлық емес математикалық
модельдерді қолдану ұңғымалардың маңында өткізгіштігі төмен немесе
жоғары сақиналы аймақтардың болуын анықтауға және олардың әлеуетті
өнімділігінің өзгеруін есептеуге мүмкіндік береді.

Бұрын гидравликалық жаруға ұшыраған ұңғымалардағы ағынды қарқындату
әдістерін енгізу жеке мәселе болып табылады, онда пайдалану процесінде
өнімділік төмендеді.

Кәсіпшілік тәжірибе көрсеткендей, жарықтардың өткізгіштігін қалпына
келтіру және ұңғымалардың өнімділігін қайта арттыру әр түрлі химиялық
композициялармен кенжар аймағын өңдеудің тиімді әдістерін енгізу арқылы
жүзеге асырылуы мүмкін. Кәдімгі ұңғымалар мен ҚГЖ бар ұңғымалардың
өнімділігін төмендету механизмі өте жақын. Сондықтан, Юра қабаттарындағы
кәдімгі ұңғымаларды өңдеу кезінде жақсы көрсеткіштерге ие болатын
жарықтардың өткізгіштігін қалпына келтіру кезінде қышқыл және басқа
химиялық өңдеу технологиясын қолданудың жоғары тиімділігін күту керек.
Осылайша, ҚГЖ-дан кейін пайдалану процесінде оны төмендеткен
ұңғымалардың өнімділігін арттыру мәселелерін Юра қабаттарында шешу
пайдалану қорының жұмысын егжей-тегжейлі талдауды және кенжар маңындағы
аймақтың өткізгіштігін төмендетудің негізгі факторларын анықтауды талап
етеді.

Жақында олар гидравликалық жару үшін ең тиімдісі тұтқырлығы төмен, әлсіз
сүзілген сұйықтық болуы керек деп санайды, өйткені бұл жағдайда қабат
бөгде заттармен бітеліп қалмайды, ал сұйықтықты құбырлар арқылы айдау
кезінде қысымның жоғалуы аз, ал ұңғыманың кенжарындағы қысым тез артады
және жарықшақ пайда болғаннан кейін жоғары деңгейде сақталады.
Нәтижесінде жарықтардың көбірек ашылуы және оларды құммен жақсы толтыру
қамтамасыз етіледі. Жыртылған сұйықтықтардың сүзілуін төмендететін
қоспалар ретінде ВНИИ асфальтит пен мұнай қоспаларын қолдануды ұсынады.

Тұтқырлықтың жоғарылауы, сондай-ақ қабаттардың жарылуы кезінде
қолданылатын сұйықтықтардың сүзілгіштігінің төмендеуі оларға тиісті
қоспаларды енгізу арқылы жүзеге асырылады. Қабаттардың жарылуы кезінде
қолданылатын көмірсутекті сұйықтықтарға арналған мұндай қоюландырғыштар
органикалық қышқылдардың тұздары, мұнайдың жоғары молекулалық және
коллоидты қосылыстары (мысалы, мұнай гудроны және басқа да мұнай өңдеу
қалдықтары) болып табылады.

Карбонатты коллекторлардың жарылуы кезінде қолданылатын кейбір мұнайлар,
керосин қышқылдары және мұнай қышқыл эмульсиялары және су-мұнай
эмульсиялары айтарлықтай тұтқырлыққа және құмның жоғары сыйымдылығына
ие. Бұл сұйықтықтар мұнай ұңғымаларында жарылған кезде жарылғыш
сұйықтықтар мен құм тасымалдаушы сұйықтықтар ретінде қолданылады.

Жүргізілген есептеулер ұңғыманың маңында өткізгіштігі төмендеген аймақ
болған кезде оның өнімділігі бірнеше рет төмендейтінін көрсетеді. Бұл
жағдайда қабаттың зақымдану тереңдігін ескере отырып жасалған табысты
қышқылдық өңдеу және басқа да кенжар аймағына әсер ету әдістері
ұңғымалардың өнімділігінің едәуір артуына әкеледі. Ұңғымалардың
шығымдарының 2 есе артуы, әсіресе өнімді бөлімде өткізбейтін секіргіштер
болған кезде, қабаттың бекітілген аралықтарын қосу кезінде де тән.
Тиісінше, өңдеудің жеткілікті тереңдігін қамтамасыз етпейтін қышқыл
құрамының көлемін дұрыс есептемегенде, экспозиция объектісінің нақты
геологиялық-физикалық жағдайлары үшін қышқыл ерітіндісіне қойылатын
барлық талаптар сақталса да, бекітілген ұңғыманың өнімділігін едәуір
арттыру мүмкін емес. Керісінше, өткізгіштігі төмен аймақ немесе (және)
шектелген аралықтар болмаған жағдайда, ұңғыма шығымының өсімі базалық
деңгейдің алғашқы ондаған пайызына дейін құрайды. Мұндай жағдайларда
ұңғымаларды қарқынды өңдеу шығындары өтелмейді.

Сүзгілеу гетерогенділігінің жоғарыда аталған нұсқалары әр түрлі
ұңғымаларда бір кен орнында да, бір ұңғымада да орын алады. Кенжар
маңындағы аймақтың сүзу гетерогенділігінің әрбір түрі, сондай-ақ
аутопсия дәрежесі мен сипаты бойынша жетілмегендік сұйықтықтарды сүзу
геометриясына және ұңғымалардың өнімділігіне әсер етеді.

Модельдің схемасы келесі сүзгілеу гетерогенділігінің нұсқаларын қамтиды:

- қабатқа сырттан әсер етілмейді, ал кенжар аймағының өткізгіштігі
қабаттың өткізгіштігіне тең;

- ұңғыма оқпанының маңында терінің оң әсерімен, сондай-ақ белгілі бір
өлшемдері мен өткізгіштігімен сипатталатын нашар өткізгіштік аймағы бар;

- ұңғыма оқпанынан ең жақын аймақ сүзу сипаттамаларына ие, қабаттың сүзу
қасиеттерінен жоғары және теріс тері әсерімен, сондай-ақ осы аймақтың
белгілі бір геометриясымен сипатталады;

- ұңғыманың маңында өткізгіштігі өзгерген екі аймақтың болуы өте жиі
кездеседі, олардың біреуі қабаттың өткізгіштігіне ие, ал екіншісі төмен
сүзу сипаттамаларына ие;

- кенжар аймағының жекелеген аралықтары толығымен шектелген.

{\bfseries Нәтижелер және талқылау.} Сүзгілеу гетерогенділігінің барлық осы
түрлері, соның ішінде аутопсия дәрежесі мен сипаты бойынша жетілмегендік
(ұңғыманың маңында өткізгіштігі жақсарған аймақты қоспағанда) қосымша
сүзу кедергісін тудырады, сондықтан оның өнімділігін төмендетеді. Кенжар
маңы аймағының жай-күйін бағалаудың және қабаттың сүзу параметрлерін
анықтаудың ең сенімді құралы қысымды қалпына келтіру қисығын жаза
отырып, ұңғымаларға стационарлық емес режимде гидродинамикалық
зерттеулер жүргізу болып табылады. Қабатты-гетерогенді қабаттарда
өткізгіштігі нашарлаған немесе толық колматациямен қатар, кенжар
аймағының өткізгіштігі қабаттың өткізгіштігінен жоғары болатын аралықтар
болған кезде күрделі жағдайлар байқалады. Мұндай жағдайларда тек
гидродинамикалық зерттеулер жүргізу кенжар маңындағы аймақтың сүзу
кедергісі туралы толық ақпарат алуды қамтамасыз етпейді. Мұны істеу үшін
оның жұмыс режиміне сәйкес депрессия кезінде ұңғымаға сұйықтық ағынының
профилін одан әрі анықтау қажет. Бұдан әрі Юра кен орындарын игерудің
ағымдағы сатысында ұңғымаларға гидродинамикалық және
кәсіптік-геофизикалық зерттеулер жүргізу әдістері мен схемалары
қаралатын болады.

Құрамында қышқыл бар қосылыстармен сазқышқылды өңдеу (СҚӨ),
сазқышқылды-тұзқышқылды өңдеу (СТҚӨ) қабаттың кенжар маңы аймағын өңдеу
тау жыныстарында қышқылмен еритін карбонатты материалдарды шаймалау
салдарынан қабаттың кенжар маңы аймағының өткізгіштігін арттыруға және
бұрғылау кезінде сазды ерітіндімен бітелген тесіктерді тазартуға
мүмкіндік береді; ұңғыманың кенжарының сүзгі бетін коррозиялық
шөгінділерден және тұндырғыш материалдардан тазартуға; коррозиялық
материалдарды ерітуге және бетіне шығаруға сорғы-компрессорлық құбырлар
(СКҚ) қабырғаларынан шөгінділер және қоршау бағаналары.

Мұнай -- газ өндіру басқрмасы (МГӨБ) кен орындары жағдайында мұнай
өндіруді арттыру мен мұнай өндіруді қарқындатудың қышқылдық әдістерін
қолдану жұмыс жүргізілген кен орындарының игерудің бастапқы, сусыз
сатысында болуына байланысты өте жоғары тиімділік береді.

Қабатқа әсер ету технологиясын құру бойынша эксперименттік жұмыстар
келесі зерттеу түрлерін қамтыды:

- коллектордың минералогиялық құрамын анықтау;

- тау жынысындағы саз материалының түрі мен мөлшерін диагностикалау;

- коллектордың литологиялық құрамын, құрылымы мен құрылымын
петрографиялық зерттеу;

- тау жыныстарының сүзу сипаттамаларын анықтау;

- тау жыныстарының физикалық-механикалық сипаттамаларын бағалау;

- реакция өнімдерін қабаттан депрессиялық алудың оңтайлы жағдайларын
анықтау.

Қабаттың кенжар маңы аймағының жай-күйі оның диагностикасымен
анықталады, оның ішінде ұңғыманың тарихын зерттеу, кен орны бойынша
өкілді негізгі материалдарды сипаттау бойынша тау жыныстарының заттық
және минералогиялық құрамдарын бағалау және гидродинамикалық
зерттеулер жүргізу. Гидродинамикалық зерттеулердің нәтижелері бойынша
реагенттерді айдау көлемі анықталады. Зерттеу деректерін (мысалы, ҚҚКҚ
(Қысымды Қалпына Келтіру Қисығы)) белгілі әдістеме бойынша түсіндіру
арқылы {[}11{]} өзгертілген өткізгіштік \(r_a\) және белсенді дренаж
радиустары \(r_a\) келесі формулалар бойынша есептеледі:

\begin{equation*}
r_n=(1.5\div 2)\sqrt{\chi_1T_1},
\end{equation*}
\vspace{-0.5cm}
\begin{equation*}
r_a=(1.5\div 2)\sqrt{\chi_2T_2},
\end{equation*}

мұндағы \(\chi_1,\chi_2\) - ҚҚКҚ - ның бастапқы және аралық учаскесі
үшін тиісінше айқындалатын қабаттың пьезоөткізгіштігі (төмен өткізгіш
коллекторы бар қабаттың кенжар аймағы үшін ҚҚКҚ \(S\) -тәрізді сипатқа
ие); \(T_1\) - ҚҚКҚ - ның бастапқы учаскесінен аралық учаскесіне өту
кезінде иілу нүктесіне сәйкес келетін уақыт; - \(T_2\) аралық
учаскеден соңғы учаскеге өту кезінде иілу нүктесіне сәйкес келетін
уақыт (қысымды толық қалпына келтіру аймағы) ҚҚКҚ учаскесі.

Химиялық реагенттерді айдау көлемі формула бойынша анықталады:

\begin{equation*}
V=\pi mh(r_{x^2}-r^2),
\end{equation*}

мұндағы \(r_x\) - мән \(r_n\) немесе \(r_a\);

\(m\) - кеуектілік коэффициенті;

\(h\) - қабаттың тиімді перфорацияланған қалыңдығы.

Орташа әсер ететін реагенттерді айдау көлемі, ол үшін \(r_x=r_a\)
қышқыл құрамын қоспағанда, \(r_x=r_a\) параметрді пайдалана отырып
айқындалады. Есептелген көлем сонымен қатар суды қайта сату көлемін
қамтиды, ал қышқыл құрамы мен су көлемінің арақатынасы 1:1,5: 2,5
құрайды.

Осылайша, кешенді химиялық депрессиялық әсерді қолдана отырып,
өткізгіштігі төмен коллекторлары бар ұңғымалардың шұңқырлы аймағын терең
өңдеу технологиясы перспективті болып табылады және гидродинамикалық
байланысқан айдау және өндіру ұңғымаларын қамтитын игеру элементіне әсер
ету кезінде қорларды өндіруді едәуір күшейтуге мүмкіндік береді.

Сынау агентінің концентрациясы, жарықшақты жабу (жабу) кернеуі және
ұңғыманың өнімділігі арасындағы тәуелділіктер қолданылады. Жарықшақты
жабу кернеуі жарықшақ бойымен тұрақты болуы керек және ұңғыманы өңдеу
(қарқындату) кезіндегі қысым мен оны іске қосу кезіндегі ұңғымадағы
қысым арасындағы айырмашылық ретінде есептелуі керек. Өткізгіштікпен
қатар модельдеу кезінде ең маңызды параметр-жарықшақтың биіктігі.
Геологиялық-каротаждық жұмыстар (ГКЖ) нәтижелері әктастар үшін тұз,
ацетон және құмырсқа қышқылдарын қолдану арқылы алынады. Бұл модельдің
көмегімен қышқылдың ену ұзындығы мен жарықшақтың өткізгіштігін есептеуге
болады. Осы модельге сәйкес ГКЖ екі кезеңде жүзеге асырылады:
біріншісінде - сынғыш агентті енгізу (үзіліс сұйықтығы), екіншісінде -
қышқыл айдау. Жарықшақтың белгілі биіктігінде жарықшақтағы қышқыл
ағынының өрісін орнатуға болады, оның ену ұзындығын анықтау үшін
жарықшақтың төменгі жағындағы қышқылдың реакция жылдамдығын білу қажет.

Зерттеу барысында доломиттерде кинетикалық реакция әктастарға қарағанда
мүлдем басқаша жүреді. Қышқыл концентрациясын қышқылдың жарықшаққа ену
қашықтығымен байланыстыратын шешім алынды. Бұл жағдайда қышқылдың
жоғалуы болмайды деп болжанады, қышқылдың жалпы концентрациясының оның
айдау аймағындағы қышқыл концентрациясына қатынасы, біреуіне тең,
қышқылдың енуінің мүмкін болатын максималды қашықтығын көрсетеді және
сәйкесінше ГКЖ кезінде жарықшақтың ұзындығын анықтайды. Модельдерде тері
әсерінің екі түрі қабылданды: ұңғыманың кенжар аймағының бұзылуына
байланысты сүзу параметрлерінің нашарлауына байланысты ламинарлы ағынның
тері әсері және сұйықтықтың Дарси заңынан ауытқуына байланысты тері
әсері.

Бұл жағдайда бірінші, ламинарлы тері әсері тұрақты болып қалады, ал
екіншісі ағынның жылдамдығына (шығымына) байланысты.

Қабаттың өткізгіштігінің төмендеуінің алдын алу жөніндегі бірінші
кезектегі іс-шаралар қабаттың ең аз зақымдануымен тиімді кептелу
сұйықтықтарын қолдану болуы тиіс. Юра шөгінділеріндегі кенжар аймағының
зақымдануының негізгі факторларын қарастыру ұңғымалардың өнімділігін
қалпына келтіру және арттыру үшін мамандандырылған көп компонентті
қышқыл ерітінділерін әзірлеу және қолдану қажеттілігі туралы қорытындыға
әкеледі. Олардың құрамында коррозия ингибиторлары және темір мен басқа
да тұздардың тиімді тұрақтандырғыштары, эмульсиялардың пайда болуын және
жойылуын болдырмауға арналған деэмульгаторлар, сондай-ақ қышқыл
ерітіндісін бейтараптандыруға дейін және одан кейін мұнаймен шекарадағы
фазааралық керілуді төмендетуге арналған арнайы қоспалар болуы керек.
Қабаттың сумен және эмульсиямен бітелуі жағдайында деэмульгаторлар мен
гидрофобизаторларды өз бетінше қолдану тиімді әдіс болады. Коллектордың
қышқыл ерітінділерімен қайталама колматациялануын бір мезгілде алдын ала
отырып, қабаттың әлеуетті өнімділігін қалпына келтіру жөніндегі міндетті
технологиялық іс-шараларға терең енетін перфорация әдістері, қабаттың
қалыңдығының 1 метріне 5-8 м\tsp{3} дейін жұмыс
ерітіндісінің үлестік көлемімен көп көлемді терең қышқылды өңдеу,
сондай-ақ қышқылмен өңдеу алдында СКҚ орналасуын ойып алу жолымен
кеуекті арналардың темір тұздарымен ластануының алдын алу жатады. Төмен
өткізгіш коллекторлардағы көп көлемді қышқылды өңдеудің болашағы мен
жоғары тиімділігі.

Кен орындарын игерудің тиімділігі, ең алдымен, оны ашу процесінде де,
ұңғымаларды пайдалану кезінде де әртүрлі физика-химиялық өзгерістерге
ұшырайтын қабаттың кенжар маңы аймағының (ҚКМА) жай-күйімен анықталады.
Сондықтан сүзу қасиеттерін жақсарту, қабылдау ағынының профильдерін
теңестіру, суландыру конустары мен бағаналы ағындарды жою бойынша
іс-шаралар жүргізу қажет.

ҚГЖ технологиясы қабатта төмен өткізгіш коллектордан мұнай алуды едәуір
күшейтуге мүмкіндік беретін төмен сүзуге төзімділігі (ТСТ) арналар
жүйесін құрумен сипатталады. Бұл процесс ҚКМА -дағы ҚГЖ -дан ұзағырақ
жарықтардың пайда болуымен, сондай-ақ оларды бұрғылаудан шығару кезінде
төмен өткізгіш қабатты ашатын барлық ұңғымаларда жүргізілуімен
ерекшеленеді.

ҚГЖ диапазонындағы сұйықтықтың (мұнайдың) сүзілу заңдылығын зерттеу үшін
біз жарылған кеуекті қабаттардың моделін қолданамыз, мұнда тау жынысы
қатты біртекті орта ретінде қарастырылады, оның кез-келген нүктесінде
қос кеуектілік, өткізгіштік және сүзу жылдамдығы болады (жарықтар мен
кеуектер жүйесі үшін). Уақыттың бастапқы сәтінде тесіктер мен
жарықтардағы қысым нөлге тең болады деп санаймыз.

\(Q_0\) - Тұрақты ағынмен нөлдік радиустық ұңғымалар арқылы сұйықтықты қабатқа
айдау кезінде модельдік есепті қарастырамыз. Содан кейін жарықтардағы
қысымға қатысты шешім келесідей болады:

\begin{equation}
p=\frac{Q_0}{4\pi\varepsilon}\left[\ln\frac{2.25kt}{r^2}-Ei(-\frac{t}{\tau_0})\right]
\end{equation}

мұндағы \(t\) - уақыт, \(c\); \(r\) - ұңғыма осінен қарастырылып
отырған нүктеге дейінгі қашықтық, \(м\); \(\varepsilon\) - жарықтар
жүйесінің гидроөткізгіштік коэффициенті, \(м^3/Па\cdot с\); \(к\) -
ортаның пьезоөткізгіштігі, \(м^2/с\); \(\tau_0=\frac{\eta}{к}\) -
кешігу уақыты, \(с\);\(\eta\) - ортаның жарылуын сипаттайтын параметр,
(жарықтар болмаған кезде \(\eta=0\)); \(Ei(z)\) - интегралды
көрсеткіштік функция.

Ортаның ерікті нүктесіндегі қысымның уақытқа тәуелділігін талдау қысым
секірісінің шамасы ұңғымада және қабат нүктелерінде уақыт аралығында
сақталатынын көрсетеді

\begin{equation*}
\frac{r^2}{к}<<t<<\tau_0
\end{equation*}

Ағынның секіргіш өзгеруінен кейін ұңғымадағы қысымның мұндай секіргіш
өзгеруі жарылған кеуекті ортаның сипатымен түсіндіріледі: бұзылыстың
таралуы жарықтар мен кеуекті блоктардағы қысымдар теңестірілгеннен
кейін, яғни уақыт өткеннен кейін пайда болады \(t\approx\tau_0\).
Мұндай тұжырым ҚГЖ жүргізгеннен кейін қысымды теңестіру уақытын анықтау
үшін маңызды {[}6{]}.

Батыс Қазақстан кен орындарында мұнай өндіруді тұрақтандыру немесе
ұлғайту мынадай жағдайлармен күрделене түсті:

1) кен орындарының игерудің кеш сатысына енуі;

2) қорлар құрылымының күрт нашарлауы (алу қиын қорлардың үлесі артады);

3) пайдалануға жеке көзқарасты талап ететін құрылымы бойынша әртүрлі кен
орындарының көп санының болуы;

4) негізгі құралдарды жаңғырту үшін қаржы ресурстарының жетіспеушілігі.

Іздеу аумағы уақытша жабылғандықтан, жаңа кен орындарының ашылуы
қауіпті, бұл айтарлықтай уақыт кезеңінен кейін әсер етуі мүмкін.

Сондай-ақ, мұнай өндіруді қарқындатудың және қабаттардың мұнай
өндірісін ұлғайтудың прогрессивті технологияларының салалық ауқымында
есепке алуды жүргізген жөн, мұнда қолданылатын әдіс бойынша бір ұңғыма
операциясы үшін инновациялық көрсеткіш берілген. Қолданылатын әсер ету
түрінің тиімділігін бағалау және \(k_{iu}\) - басымдылық көрсеткішін
есептеу үшін бір ұңғыма операциясы үшін формула ұсынылады:

\begin{equation}
k_{iu}=\frac{1}{3}\left(\frac{N_{min}}{N_i}+\frac{Э_i}{Э_{max}}+\frac{З_{min}}{З_i}\right)
\end{equation}

мұндағы \(Э_i\), \(N_i\) және \(З_i\) - сәйкесінше, қосымша қосымша
өндіру, ұңғымалық операциялар саны және шығындар (мың
доллар.АҚШ) \(i\)-ші әсер ету түрі үшін
\((i=1,2,3\ldots9)\) \(N_{min}\), \(Э_{min}\) және \(З_{min}\) -
сәйкесінше ұңғымалық операциялардың минималды саны, максималды мұнай
өндіру және барлық әсер ету түрлерінің ең аз шығыны (әсер етудің
үшінші түрі алынып тасталады).

Барлық геологиялық-техникалық іс-шаралар (ГТШ) екі үлкен топқа бөлінеді:
жылу, газ және физика - химиялық әсерлерді қамтитын қабаттардың мұнай
бергіштінін арттыру әдістері (МБАӘ); қабатқа әсер ету әдістері және
ұңғыма технологиялары-ҚГЖ, жүйенің жұмыс режимдерін реттеу
технологиялары.

Бірінші топтың әдістері дәстүрлі су айдау әдісімен салыстырғанда
қорларды өндіруді жақсартуға бағытталған принципті технологиялық
шешімдер кешені болып табылады. Бұл әдістердің айрықша ерекшелігі-әрбір
нақты жағдайда ғылыми-зерттеу жұмыстарын жүргізу үшін ғылыми --
техникалық прогрессті қолдану қажеттілігі, сондай-ақ шығындар сипаты,
технологиялық және экономикалық тәуекелдің жоғарылауы.

Екінші топтың әдістері қазіргі заманғы ғылым мен практикада сәтті
сыналған әдістер мен технологиялар болып табылады, олар әзірленген
бағдарламалық-имитациялық модельдер негізінде типтік тапсырмалар мен
шешімдерді қолданады. Осы топ үшін қабаттың мұнай бергіштігін арттыру
(ҚМБА) әдістерін енгізу тәжірибесі жоғары экономикалық көрсеткіштері бар
ең тиімді технологияларды таңдауға мүмкіндік береді.

Мұнай бергіштік коэффициентін (МБК) арттыру әдісін қолданудың
техникалық-экономикалық негіздемесі кезінде бір кен орны шегінде
ұңғымалар тобы бойынша және геологиялық-физикалық сипаттамалары жақын
игеру ауданы шегінде кен орындары тобы бойынша әдістің тиімділігін
белгілеу басты рөл атқарады. Осы мақсат үшін әдістің жарамдылық
дәрежесі ұңғымаларды пайдалануға беру және осы әдісті қолдана отырып
кен орындарын игеру кезектілігін анықтайтыны анық. Әдістің жарамдылық
дәрежесін бағалау қарастырылып отырған ұңғыманы немесе кен орындарын
неғұрлым толық сипаттайтын және таңдау белгілерінің рөлін
атқаратын \(x_i\) параметрлер жиынтығы бойынша белгіленеді.

Қабаттардың өнімділігін арттырудың объективті мүмкіндігі қарастырылып
отырған объектілердің ҚГЖ әдісінің белгілі бір модификацияларын
қолдануға жарамдылығына қарай артады. Ұңғымалар бойынша қолданылатын
әсер ету түрінің тиімділігін таңдаудың 7 белгісі арқылы бағалаймыз:
коллекторлардың тиімді қалыңдығы \(h(м)\), объектінің тереңдігі \(L(м)\),
сұйықтықтың көлемді айдалуы \(Q^{(айдау)}(м^3)\),
сулану (пайызбен) \(Q^{(қос)}\), өндірілетін сұйықтық көлемі
\(Q^{(өнд)}\varepsilon\), қосымша мұнай \(Q^{(қос)}тонна/сут\),
ұңғыманың жұмыс уақыты \(T_{(сағ)}\).

\(і\) - ұңғыма әдісінің \(k_i\) - жарамдылық көрсеткішін формула арқылы анықтаймыз:
\end{multicols}
\vspace{-0.3cm}
\[k_{i\ } = \frac{1}{7}\left\lbrack \frac{L_{min}}{L_{i}} + \frac{h_{i}}{h_{\max}} + \frac{Q_{\min}^{(айд)}}{Q_{i}^{(айд)}} + \frac{Q_{min}^{(өнд)}}{Q_{i}^{(өнд)}} + \frac{Q_{\min}^{(сул)}}{Q_{i}^{(сул)}} + \frac{Q_{i}^{қос}}{Q_{\max}^{қос}} + \frac{T_{\min}}{T_{i}} \right\rbrack\]

\begin{multicols}{2}
Мұндағы, \(L_{\min}{,h_{\max},Q}_{\min}^{(айд)}\),
\(Q_{\min}^{(өнд)}\), \(Q_{\min}^{(cул)}\), \(Q_{\max}^{(қос)}\), \(T_{\min}\)
-арқылы тиісті үлгілердің экстремалды мәндері. Батыс Қазақстан \emph{Х}
және \emph{У} екі кен орнындағы үш ұңғыма үшін ҚГЖ әдісінің тиімділік
көрсеткішін анықтау әдісін қолданамыз (1, 2-кесте). Есептеулерде ҚГЖ
жүргізілгеннен кейін ұңғымалардың жұмыс режимінің деректері пайдаланылды
{[}6{]}.
\end{multicols}

\tcap{1-кесте. Х кен орнының геологиялық-физикалық параметрлері}
\begin{longtblr}[
  label = none,
  entry = none,
]{
  width = \linewidth,
  colspec = {Q[20]Q[462]Q[143]Q[83]Q[83]Q[83]},
  cells = {c},
  column{2} = {l},
  cells = {font = \small},
  cell{1}{1} = {r=2}{c, font=\bfseries\small},
  cell{1}{2} = {r=2}{c, font=\bfseries\small},
  cell{1}{3} = {r=2}{c, font=\bfseries\small},
  cell{1}{4} = {c=3}{c, 0.249\linewidth,font=\bfseries\small},
  hlines,
  vlines,
}
№ & Параметрлер                                           & Экстремалды & Ұңғымалар &       &       \\
  &                                                       &             & №654      & №525  & №655  \\
1 & Тиімді қалыңдық, \(h(м)\)                             & 7           & 6         & 6     & 7     \\
2 & Жатыс тереңдігі, \(L(м)\)                             & 2251        & 2850      & 2742  & 2251  \\
3 & Айдау көлемі, \(Q^{(айд)}(м^3)\)                       & 47          & 47        & 81    & 85    \\
4 & Өндірілген сұйықтық көлемі, \(Q^{(сул)}тонн/тәулік\)   & 3           & 22,2      & 25,4  & 3     \\
5 & Сулану, \(Q^{(сул)}\)(\%)                              & 5           & 8         & 8     & 5     \\
6 & Қосымша мұнай көлемі, \(Q^{(қос)}\) \(тонн/тәулік\)     & 8241        & 104       & 8241  & 1067  \\
7 & Жұмыс сағаттарының саны \(T(10^3 \text{сағат})\)      & 10,08       & 10,08     & 11,49 & 10,99 \\
8 & Әдістің тиімділік көрсеткіші                          & -           & 0,635     & 0,697 & 0,800 
\end{longtblr}

\tcap{2-кесте. У - кен орнының геологиялық-физикалық параметрлері}
\begin{longtblr}[
  label = none,
  entry = none,
]{
  width = \linewidth,
  colspec = {Q[20]Q[462]Q[143]Q[83]Q[83]Q[83]},
  cells = {c},
  column{2} = {l},
  cells = {font = \small},
  cell{1}{1} = {r=2}{c, font=\bfseries\small},
  cell{1}{2} = {r=2}{c, font=\bfseries\small},
  cell{1}{3} = {r=2}{c, font=\bfseries\small},
  cell{1}{4} = {c=3}{c, 0.249\linewidth,font=\bfseries\small},
  hlines,
  vlines,
}
№ & Параметрлер                                           & Экстремалды & Ұңғымалар &       &       \\
  &                                                       &             & №84       & №88   & №69   \\
1 & Тиімді қалыңдық, \(h(м)\)                             & 7           & 6         & 6     & 7     \\
2 & Жатыс тереңдігі, \(L(м)\)                             & 2255        & 2254      & 2255  & 2255  \\
3 & Айдау көлемі, \(Q^{(айд)}(м^3)\)                       & 55          & 55        & 83    & 83    \\
4 & Өндірілген сұйықтық көлемі, \(Q^{(өнд)}тонн/тәулік\)   & 5,6         & 32,5      & 11,7  & 5,6   \\
5 & Сулану, \(Q^{(сул)}\) (пайыз)                          & 20          & 20        & 75    & 87    \\
6 & Қосымша мұнай көлемі, \(Q^{(қос)}\) \(тонн/тәулік\)    & 188,8       & 188,8     & 79,1  & 66,3  \\
7 & Жұмыс сағаттарының саны \(T(10^3 \text{сағат})\)      & 10,03       & 11,06     & 10,48 & 10,03 \\
8 & Әдістің тиімділік көрсеткіші                          & -           & 0,848     & 0,663 & 0,750 
\end{longtblr}

\begin{multicols}{2}
Кестелік мәліметтерден Х кен орны үшін тиімділіктің ең үлкен көрсеткіші
№655 ұңғымаға ие екендігі шығады, өйткені мұнда өнімді қабаттың тиімді
қалыңдығы максималды мәнге ие және құрамында қосымша мұнай бар
өндірілетін сұйықтық көлемі ең аз мәнге ие. Сонымен қатар, бұл ұңғыманың
тереңдігі мен сулануының ең кіші мәндері бар. Бұл үлгілердің мәндері
ұңғыманың өнімділігін арттыру үшін әдістің жарамдылық дәрежесін
түпкілікті бағалау үшін оң рөл атқарады. Кен орын үшін № 84 ұңғыма ең
жоғары көрсеткішке ие, мұнда қосымша мұнай өндіру және төмен сулану
шешуші рөл атқарады.

{\bfseries Қорытынды.} Іздестіру-барлау жұмыстарын жүргізуге уақытша шектеу
жағдайында жаңа кен орындарын ашу перспективалары өте белгісіз болып
қалады. Мұндай жағдайларда геологиялық барлау жоғары тәуекелдермен
байланысты және сәтті болса да, ұзақ уақыт өткеннен кейін ғана нақты
нәтиже бере алады. Сондықтан Батыс Қазақстанның жұмыс істеп тұрған мұнай
объектілерінде бүгінде өндіру көлемін ұлғайтудың бірден --- бір
мүмкіндігі-игерілген қабаттардың мұнай беруін арттыруға және
қолданыстағы кен орындарын игеруді оңтайландыруға бағытталған
қарқындатудың заманауи технологияларын-әдістерді енгізу болып табылады.

Батыс Қазақстанның бірнеше кен орындарын игеру объектілері үшін
таңдаудың 8 белгісі арқылы ұңғымалар бойынша қолданылатын әсер ету
түрінің тиімділігі бағаланды. Өндіру деңгейін арттырудың объективті
мүмкіндігі қарастырылып отырған объектілердің ҚГЖ әдісінің белгілі бір
модификацияларын қолдануға жарамдылығына қарай артады.

ҚГЖ-дан кейінгі көрсеткіштерді нашарлатқан мұнай ұңғымаларының шығымын
арттыру мәселесін шешу қабатқа әсер етудің қосымша әдістерін қолдануды
талап етеді. Юра қабаттарында пайдалану қорының жұмысын егжей-тегжейлі
талдау, сондай-ақ кенжар аймағының өткізгіштігінің төмендеуінің негізгі
факторларын анықтау қажет.
\end{multicols}

\begin{center}
{\bfseries Әдебиеттер}
\end{center}

\begin{refs}
1. Батыров М.И. Техника и технология проведения ГРП на Ем-Ёговской
площади Красноленинского месторождения // Отраслевые научные и
прикладные исследования: Науки о Земле. -2020. -№ 4. -С.89--106.

2. Султанов Ш.Х., Мухаметшин В.Ш., Стабинскас А.П. и др. Исследование
возможности использования воды с высокой минерализацией для
гидравлического разрыва пласта // Записки Горного института. -2024.
-Т.270. - С.950-962.

3. Казинский А.А. и др. Требования к жидкостям разрыва при проведении ГРП
на месторождении Карачаганак //
\href{https://ojs.wkau.kz/index.php/gbj/issue/view/79}{Наука и
образование}. -2023. -T.2. -№ 4 (73). -С.265-273. DOI
\href{https://doi.org/10.52578/2305-9397-2023-4-2-265-273}{10.52578/2305-9397-2023-4-2-265-273}.

4. Галлямов А.Р. Эффективность применения гидроразрыва пласта //
Международный научный журнал «Вестник Науки». -2025. --T.5. -№ 8(89).
-С.327-334.

5. Исламов Я.Р., Нурлыгаянова Э.Р., Вагизов А.М. и др. Развитие
технологий гидравлического разрыва пласта для карбонатных отложений
каширского и подольского горизонтов Арланского нефтяного месторождения
// Георесурсы. - 2024. -Т.26(4). -С.200-208. DOI
10.18599/grs.2024.4.21.

6. Бухарбаева А.Н., Асанов К.Б., Башев А.А., Джаксылыков Т.С., Марданов
А.С. Перспективы повышения эффективности гидроразрыва пласта с
использованием высокотехнологичных проппантов на месторождениях
Атырауского региона// Вестник нефтегазовой отрасли Казахстана. -2024.
-Т.6(4). -C.49-67. DOI
\href{https://doi.org/10.54859/kjogi108775}{10.54859/kjogi108775}.

7. Асанов К.Б. Опыт применения многостадийного гидроразрыва пласта в
обсаженном стволе вертикальной скважины. // Вестник нефтегазовой
отрасли Казахстана. -2019. -№ 1. -C.88-94. ISSN 2707-4226.

8. Волокова М.А., Петухова Е.В. Гидравлический разрыв пласта
высокоэффективный метод интенсификаций добычи нефти //Успехи в химии и
химической технологии. -2023. --T.37(4). - C.52-63.

9. Бисенгалиев М.Д., Баямирова Р.У., Жолбасарова А.Т., Тогашева А.Р.,
Сарбопеева М.Д. Мероприятия по улучшению эффективности гидравлического
разрыва пласта и Скин-ГРП // Журнал нефти и газ. -2023. -T.1 (133). -
C.81-93. DOI
\href{https://doi.org/10.37878/2708-0080/2023-1.07}{10.37878/2708--0080/2023-1.07}.

10. Мазо А.Б., Хамидуллин М.Р. Явно-неявные алгоритмы ускорения расчета
двухфазного притока к горизонтальной скважине с многостадийным
гидроразрывом пласта// Вычислительные методы и программирование.
-2017. -Т.18(3). - С.204 -- 213. DOI 10.26089/NumMet.v18r318.

11. Mukherjee H., Economides M.J. A parametric comparison of horizontal
and vertical well performance // SPE Form Eval. -1991. -Vol.6 (2).
-P.209-216. DOI 10.2118/18303- PA.
\end{refs}

\begin{center}
{\bfseries References}
\end{center}

\begin{refs}
1. Batyrov M.I. Tehnika i tehnologija provedenija GRP na Em-Jogovskoj
ploshhadi Krasnoleninskogo mestorozhdenija // Otraslevye nauchnye i
prikladnye issledovanija: Nauki o Zemle. -2020. -№ 4. -S.89-106. {[}in
Russian{]}

2. Sultanov Sh.H., Muhametshin V.Sh., Stabinskas A.P. i dr. Issledovanie
vozmozhnosti ispol' zovanija vody s vysokoj
mineralizaciej dlja gidravlicheskogo razryva plasta // Zapiski Gornogo
instituta. -2024. -T.270. - S.950-962. {[}in Russian{]}

3. Kazinskij A.A. i dr. Trebovanija k zhidkostjam razryva pri provedenii
GRP na mestorozhdenii Karachaganak // Nauka i obrazovanie. -2023. -T.2.
-№ 4 (73). -S.265-273. DOI 10.52578/2305-9397-2023-4-2-265-273. {[}in
Russian{]}

4. Galljamov A.R. Jeffektivnost'{} primenenija
gidrorazryva plasta // Mezhdunarodnyj nauchnyj zhurnal «Vestnik Nauki».
-2025. --T.5. -№ 8(89). -S.327-334. {[}in Russian{]}

5. Islamov Ja.R., Nurlygajanova Je.R., Vagizov A.M. i dr. Razvitie
tehnologij gidravlicheskogo razryva plasta dlja karbonatnyh otlozhenij
kashirskogo i podol' skogo gorizontov Arlanskogo
neftjanogo mestorozhdenija // Georesursy. - 2024. -T.26(4). -S.
200-208. DOI 10.18599/grs.2024.4.21. {[}in Russian{]}

6. Buharbaeva A.N., Asanov K.B., Bashev A.A., Dzhaksylykov T.S.,
Mardanov A.S. Perspektivy povyshenija jeffektivnosti gidrorazryva plasta
s ispol' zovaniem vysokotehnologichnyh proppantov na
mestorozhdenijah Atyrauskogo regiona// Vestnik neftegazovoj otrasli
Kazahstana. -2024. -T.6(4). - C.49-67. DOI 10.54859/kjogi108775. {[}in
Russian{]}

7. Asanov K.B. Opyt primenenija mnogostadijnogo gidrorazryva plasta v
obsazhennom stvole vertikal' noj skvazhiny. // Vestnik
neftegazovoj otrasli Kazahstana. -2019. -№ 1. -C.88--94. ISSN 2707-
4226. {[}in Russian{]}

8. Volokova M.A., Petuhova E.V. Gidravlicheskij razryv plasta
vysokojeffektivnyj metod intensifikacij dobychi nefti //Uspehi v himii i
himicheskoj tehnologii. -2023. --T.37(4). - C.52-63. {[}in Russian{]}

9. Bisengaliev M.D., Bajamirova R.U., Zholbasarova A.T., Togasheva A.R.,
Sarbopeeva M.D. Meroprijatija po uluchsheniju jeffektivnosti
gidravlicheskogo razryva plasta i Skin-GRP // Zhurnal nefti i gaz.
-2023. -T.1 (133). -C.81-93. DOI 10.37878/2708--0080/2023-1.07. {[}in
Russian{]}

10. Mazo A.B., Hamidullin M.R. Javno-nejavnye algoritmy uskorenija
rascheta dvuhfaznogo pritoka k gorizontal' noj skvazhine
s mnogostadijnym gidrorazryvom plasta// Vychislitel' nye
metody i programmirovanie. -2017.-T.18(3). - S.204 -- 213. DOI
10.26089/NumMet.v18r318.

11. Mukherjee H., Economides M.J. A parametric comparison of horizontal
and vertical well performance // SPE Form Eval. -1991. -Vol.6 (2). -P.
209--216. DOI 10.2118/18303- PA.
\end{refs}

\begin{info}
\hspace{1em}\emph{{\bfseries Авторлар туралы мәлметтер}}

Бисенгалиев М.Д. - техника ғылымының кандидаты, профессор, «С.Өтебаев
атындағы Атырау мұнай және газ университеті» КеАҚ, Атырау, Казахстан,
e-mail: maks\_bisengali@mail.ru;

Сулейменова Р.Т. - PhD, техника ғылымының кандидаты, қауымдастырылған
профессор,

«С.Өтебаев атындағы Атырау мұнай және газ университеті» КеАҚ, Атырау,
Казахстан, e-mail: raikhan.suleimenova@aogu.edu.kz;

Мухаммад Насир - PhD, Химия ғылымдары мектебі, Малайзия зерттеу
университеті, Пинанг, Малайзия, e-mail: mnm@ysm.my;

Досказиева Г.Ш. - техника ғылымының кандидаты, профессор, «С.Өтебаев
атындағы Атырау мұнай және газ университеті» КеАҚ, Атырау, Казахстан,
e-mail: doskaziyeva.gulsin@gmail.com;

Зайдемова Ж.К. - техника ғылымының кандидаты, профессор, «С.Өтебаев
атындағы Атырау мұнай және газ

Тулегенова О.Ш, - техника ғылымының кандидаты, қауымдастырылған
профессор, «С.Өтебаев атындағы Атырау мұнай және газ университеті»,
Атырау, Казахстан;

Шаяхметова Ж.Б. - техника ғылымының кандидаты, қауымдастырылған
профессор, «С.Өтебаев атындағы Атырау мұнай және газ университеті»
КеАҚ, Атырау, Казахстан, e-mail: zhanar6688@mail.ru;

Каримова А.С. - техника ғылымының кандидаты, қауымдастырылған
профессор, «С.Өтебаев атындағы Атырау мұнай және газ университеті»
КеАҚ, Атырау, Казахстан, akmaral0167@mail.ru.

\hspace{1em}\emph{{\bfseries Information about the authors}}

Bissengaliyev M.D. - Ph, professor, S. Utebaev Atyrau University of
Oil and Gas, Atyrau, Kazakhstan, e-mail: maks\_bisengali@mail.ru;

Suleimenova R.Т. - PhD, ass. professor, S. Utebaev Atyrau University
of Oil and Gas, Atyrau, Kazakhstan, e-mail:
raikhan.suleimenova@aogu.edu.kz;

Nasir M.- PhD, School of Chemical Sciences, Universiti Sains Malaysia,
Pinang, Malaysia e-mail: mnm@ysm.my;

Doskazieva G.SH. - Ph, ass. Professor, S. Utebaev Atyrau University of
Oil and Gas, Atyrau, Kazakhstan, e-mail: aing-zhomart@mail.ru;

Zaidemova Zh.R. - Ph, professor, S. Utebaev Atyrau University of Oil
and Gas, Atyrau, Kazakhstan, e-mail: b.n.m.99@list.ru;

Tulegenova O.SH. - Ph, ass. professor, S. Utebaev Atyrau University of
Oil and Gas, Atyrau, Kazakhstan, e-mail: olympiada70@mail.ru;

Shayakhmetova Zh.B.- Ph, ass. professor, S. Utebaev Atyrau University
of Oil and Gas, Atyrau, Kazakhstan, e-mail: zhanar6688@mail.ru;

Karimova A.S. - Ph, ass. professor , S. Utebaev Atyrau University of
Oil and Gas, Atyrau, Kazakhstan, e-mail: akmaral0167@mail.ru.
\end{info}
